\chapter[The Fixed Normalized Generative Law]{The Fixed Normalized\\Generative Law}
\label{ch:generative}

This chapter answers how one normalized joint law of latent samples and records is built at $c_*$ and
what passive re-expression of the frames does to it. The joint is normalized for every finite record
hypergraph, cyclic ones included (\Cref{prop:gen-exact-normalization}), local normalization of
reciprocal latent conditionals supplies no joint (\Cref{prop:gen-gibbs-counterexample}), and under
kernel covariance the evidence and the bound of a transported coupling are invariant
(\Cref{prop:gen-product-evidence-invariance,cor:elbo-independent-frame-invariance}). Its evidence is
bounded in \Cref{ch:elbo}, and attention and coarse channels act on it in
\Cref{ch:selective-communication,ch:coarsegraining}.

\paragraph{Scope.}
The cross morphisms supply no generative factor, and a model with factors in both directions needs its
own normalized joint. No quotient of the parameter space by passive re-expression is formed, and
no parameter count subtracts the dimension of the vertex group, whose action can have nontrivial
stabilizers. \status{NOT-CLAIMED}

\section{Parameters, the directed latent law, and the typing prohibition}
\label{sec:gen-typing-setup}
\label{sec:gen-composition}
\label{sec:gen-typing}

\definitionheading{Parameters, structural data, and generative factors}{def:gen-parameter-split}
Fix the agent set $V=\{1,\ldots,N\}$ and the spaces and references of \Cref{ch:probability}. Structural
data $X\in\mathsf X$ contain a finite directed acyclic graph $\Gamma$ on $V$ with parent sets
$\operatorname{pa}(i)$, the interaction multigraph $\mathfrak I$ with its links, and the frames; the
numerical parameters that a learning coordinate may move are denoted $\theta$. The split between $X$
and $\theta$ is declared, not derived, and every typing statement below applies to both. For each agent
$i$ declare the Markov kernels of the first line below, its model and state factors, which are its model
and belief priors and are generative factors, not components of the agent of \Cref{def:geo-agent}; at a
root, where $\operatorname{pa}(i)=\varnothing$, the first is a proper model law and the second a bridge
from the agent's own model coordinate to its belief coordinate. The state factor is the evaluated
kernel of \Cref{def:prob-model-evaluation} with source
$\mathsf S_i^X=\prod_{j\in\operatorname{pa}(i)}\mathsf K_j$, the empty product being a singleton
$\{\ast\}$, and target $\mathsf T_i^X=\mathsf K_i$ (second line):
\begin{gather}
P_{\theta,i}^{m}(dm_i\given m_{\operatorname{pa}(i)},X),
\qquad
P_{\theta,i}^{k}(dk_i\given k_{\operatorname{pa}(i)},m_i,X),
\label{eq:gen-kernels}\\
\operatorname{ev}_i(m_i)=K^X_{i,m_i},
\qquad
K^X_{i,m_i}(k_{\operatorname{pa}(i)},B):=P_{\theta,i}^{k}(B\given k_{\operatorname{pa}(i)},m_i,X).
\label{eq:gen-evaluated-state-kernel}
\end{gather}
With $\mathsf Y_{\operatorname{pa}(i)}=\prod_{j\in\operatorname{pa}(i)}\mathsf Y_j$, the agent mechanism is the
ordered composite below, which reads the whole parent block $y_{\operatorname{pa}(i)}$ and has unit mass:
\begin{gather}
G_{\theta,i}:(\mathsf X\times\mathsf Y_{\operatorname{pa}(i)})\rightsquigarrow\mathsf Y_i,
\label{eq:gen-agent-design-mechanism}\\
G_{\theta,i}\bigl(d(k_i,m_i)\given y_{\operatorname{pa}(i)},X\bigr)
:=P_{\theta,i}^{m}(dm_i\given m_{\operatorname{pa}(i)},X)K^X_{i,m_i}(k_{\operatorname{pa}(i)},dk_i).
\label{eq:gen-agent-mechanism}
\end{gather}

\Cref{def:obs-fast-slow-blocks} divides these coordinates into fast recognition and slow model
coordinates. At a root the state factor induces, for a model law $s$ on
$\mathsf M_i$, the belief law
\begin{equation}
\widetilde\Phi^{\mathrm{gen}}_i(s)(B)=\int_{\mathsf M_i}P_{\theta,i}^{k}(B\given m,X)s(dm),
\qquad B\in\mathscr K_i,
\label{eq:gen-root-law-bridge}
\end{equation}
which can serve as the representative at $c_*$ of the cross morphism $\widetilde\Phi$ of
\Cref{def:geo-cross-morphisms} when that identification is declared; at a nonroot the evaluated kernel
is indexed by the parent belief tuple and is a kernel-valued bridge, not a map of law fibers.

\definitionheading{Finite directed latent law}{cons:gen-finite-directed-law}
Choose a topological ordering $v_1,\ldots,v_N$ of $\Gamma$. The latent law is the ordered composition of
the first line below, where the arrow denotes ordered application of kernels as in the second line, for
bounded measurable $f$; every conditioning coordinate is introduced before its kernel is applied:
\begin{gather}
\mathbb P_{\theta,V}^{\mathsf Y}(dy\given X)
:=\overrightarrow{\prod_{l=1}^{N}}
G_{\theta,v_l}\bigl(dy_{v_l}\given y_{\operatorname{pa}(v_l)},X\bigr),
\label{eq:gen-composition}\\
\int f(y)\mathbb P_{\theta,V}^{\mathsf Y}(dy\given X)
=\int G_{\theta,v_1}(dy_{v_1}\given X)\cdots
\int G_{\theta,v_N}\bigl(dy_{v_N}\given y_{\operatorname{pa}(v_N)},X\bigr)f(y).
\label{eq:gen-iterated-integral}
\end{gather}
When the factors have densities against the component references, the densities are
\begin{gather}
g_{\theta,i}(k_i,m_i\given y_{\operatorname{pa}(i)},X)
=p_{\theta,i}^{m}(m_i\given m_{\operatorname{pa}(i)},X)p_{\theta,i}^{k}(k_i\given k_{\operatorname{pa}(i)},m_i,X),
\label{eq:gen-density}\\
p_{\theta,V}^{\mathsf Y}(y\given X)=\prod_{i\in V}g_{\theta,i}(k_i,m_i\given y_{\operatorname{pa}(i)},X),
\qquad
\mathbb P_{\theta,V}^{\mathsf Y}(dy\given X)=p_{\theta,V}^{\mathsf Y}(y\given X)\nu^{\mathsf Y}(dy).
\label{eq:gen-design-density}
\end{gather}

\hypothesisheading{Typing prohibition}{req:gen-typing-prohibition}
Each factor in \eqref{eq:gen-kernels} and each record kernel is a function only of $\theta$, $X$, and its
displayed latent coordinates; no generative kernel takes as an argument the recognition law
$\mathbb Q_{V,o,X}$, one of its marginals or parameters, or the posterior $\boldsymbol\Pi_{\theta,V,o,X}$. In
the two-argument form $\Fenergy_V(\mathbb Q;\theta)$ of the free energy \citep{Neal1998}, every evaluation
holds the generative argument $(\theta,X)$ fixed while the recognition argument varies, which excludes the
composite $\mathbb Q\mapsto\Fenergy_V(\mathbb Q;\theta(\mathbb Q))$ inside one evaluation: a generative
factor that reads the recognition law being scored, such as a transition built from a live peer
recognition marginal, is a map from recognition laws to kernels rather than a factor of one joint. A
parameter may depend on the observation or on an earlier recognition law if it is held fixed while the
bound is evaluated, so slow evolution of an agent's model or prior is admissible as recognition dynamics
of the model channel under one fixed joint or as accepted M-coordinates on $\theta$
(\Cref{prop:elbo-exact-m-coordinate,def:elbo-generalized-m-acceptance,prop:elbo-evidence-monotonicity}),
each accepted value defining a new fixed joint to which the evidence identity applies afresh.

Coupling through latents and loops of inference updates under a fixed joint remain admissible
(\Cref{sec:full-graph-closing-tower}), and a family of joints indexed by recognition laws designates no
target (\Cref{prop:gen-no-distinguished-target}).

\section{Normalized record extension and the undirected alternative}
\label{sec:gen-normalization}
\label{sec:gen-undirected}

\propositionheading{Normalization by latent composition and record attachment}{prop:gen-exact-normalization}
The iterated integral \eqref{eq:gen-iterated-integral} defines a probability kernel
$X\mapsto\mathbb P_{\theta,V}^{\mathsf Y}(\cdot\given X)$. If the record coordinates are conditionally
independent given $(X,y)$, the population law is the kernel of the first line below, and the second
line holds for every finite record hypergraph, cyclic or not, and every $(\theta,X)$:
\begin{gather}
\mathbb P_{\theta,V}(dy,do\given X)
=\mathbb P_{\theta,V}^{\mathsf Y}(dy\given X)\prod_{a\in\mathfrak A}K_a(X,y_{\partial a},do_a),
\label{eq:gen-record-extension}\\
\mathbb P_{\theta,V}(\mathsf Y\times\mathsf O\given X)=1.
\label{eq:gen-normalization}
\end{gather}
If $K_a(X,y_{\partial a},do_a)=\exp\bigl(-E_{a,o}(y_{\partial a})\bigr)\nu_a(do_a)$ with a jointly measurable
negative log density $E_{a,o}$ and $\nu^{\mathsf O}=\bigotimes_{a}\nu_a$, then
\begin{equation}
p_{\theta,V}(o,y\given X)
=p_{\theta,V}^{\mathsf Y}(y\given X)\exp\Bigl(-\sum_{a\in\mathfrak A}E_{a,o}(y_{\partial a})\Bigr)
\label{eq:gen-record-density}
\end{equation}
is its density. \status{ESTABLISHED}

\noindent\emph{Proof.}
Integrating a bounded measurable function against a probability kernel preserves boundedness and
measurability in the remaining variables, so induction along the topological order defines a measurable
kernel in $X$, and $f\equiv1$ gives unit mass; finite kernel composition gives measurability of
\eqref{eq:gen-record-extension}. By Tonelli's theorem, integrate the record coordinates first, in any
order: each $K_a$ has unit mass in $o_a$, so the finite product disappears and leaves
$\mathbb P_{\theta,V}^{\mathsf Y}(dy\given X)$ whatever the overlaps of the scopes, since the records are
conditionally independent children of one latent configuration. The same computation with the
Radon--Nikodym derivatives proves \eqref{eq:gen-record-density}. \hfill$\square$

Each record kernel is owned once by the joint, while the coordinate-local free energies of \Cref{ch:elbo}
include every factor incident to a coordinate.

\paragraph{Running example.}
The agents of \Cref{sec:geo-running-example} carry the directed joint with
$\operatorname{pa}(1)=\varnothing$, $\operatorname{pa}(2)=\operatorname{pa}(3)=\{1\}$, independent
$\mathcal N(0,1)$ model coordinates, $k_1\given m_1\sim\mathcal N(0,e^{-m_1}I_3)$, and children
transported from agent $1$,
$k_2\given k_1,m_2\sim\mathcal N\bigl(\rho_b(\Theta^b_{\bar e_1})k_1,e^{-m_2}I_3\bigr)$ and
$k_3\given k_1,m_3\sim\mathcal N\bigl(\rho_b(\Theta^b_{e_3})k_1,e^{-m_3}I_3\bigr)$. Two binary records are
attached: $a_1$ with scope $\{1,2\}$ as in \eqref{eq:geo-running-record}, and $a_2$ with scope $\{2,3\}$
and success probability $\exp(-\tfrac12\lvert k_2-\rho_b(\Theta^b_{e_2})k_3\rvert^2)$. The records close the
latent tree into the triangle, and \Cref{prop:gen-exact-normalization} normalizes the joint.
At the conditional means of the children the residual read by $a_2$ is
$\rho_b(\Theta^b_{\bar e_1})k_1-\rho_b(\Theta^b_{e_2}\Theta^b_{e_3})k_1
=\rho_b(\Theta^b_{e_1})^{-1}(I_3-\rho_b(\mathrm{Hol}^b(p)))k_1$, whose length is $\sqrt3$ times that of
the component of $k_1$ orthogonal to the rotation axis of the holonomy \eqref{eq:geo-running-holonomy},
so the records see the loop holonomy through an invariant.

\definitionheading{Normalized undirected model}{def:gen-undirected-model}
For nonnegative potentials on the agents, on the edges of an undirected graph with edge set
$E^{\mathrm u}$, and on the records, put $\widetilde p_X$ and $Z_X$ as in the first two lines below.
When $0<Z_X<\infty$, a condition that can depend on the parameters, the undirected model is the law with
the density of the third line:
\begin{gather}
\widetilde p_X(o,y)
=\prod_{i\in V}\psi_i(k_i,m_i;X)
\prod_{\{i,j\}\in E^{\mathrm u}}\psi_{ij}(y_i,y_j;X)
\prod_{a\in\mathfrak A}\psi_a(o_a,y_{\partial a};X),
\label{eq:gen-gibbs-potential}\\
Z_X
=\int\widetilde p_X(o,y)\nu^{\mathsf O}(do)\nu^{\mathsf Y}(dy),
\label{eq:gen-gibbs-partition}\\
p_{\theta,V}^{\mathrm{Gibbs}}(o,y\given X)=Z_X^{-1}\widetilde p_X(o,y).
\label{eq:gen-gibbs-density}
\end{gather}

\propositionheading{Local normalization is insufficient}{prop:gen-gibbs-counterexample}
Normalized agent potentials and an everywhere positive, finite, continuous edge potential need not give a
finite $Z_X$. \status{ESTABLISHED}

\noindent\emph{Proof.}
Take two agents with scalar belief coordinates $k_1,k_2\in\R$ under Lebesgue measure, singleton model and
record spaces with unit point mass, and potentials with product exponent
$-\tfrac12k^\top\mathsf W_{\mathrm c}k$,
\begin{gather}
\psi_i(k_i,m_i;X)=(2\pi)^{-1/2}e^{-k_i^2/2},
\qquad
\psi_{12}(y_1,y_2;X)=e^{\mathrm ck_1k_2},
\qquad \mathrm c\geq1,
\label{eq:gen-gibbs-counterexample}\\
\mathsf W_{\mathrm c}=\begin{psmallmatrix}1&-\mathrm c\\-\mathrm c&1\end{psmallmatrix},
\qquad
\operatorname{spec}\mathsf W_{\mathrm c}=\{1-\mathrm c,1+\mathrm c\};
\label{eq:gen-gibbs-counterexample-matrix}
\end{gather}
in orthonormal sum and difference coordinates the integral factorizes, and the factor in the direction of
the eigenvalue $1-\mathrm c\leq0$ diverges. \hfill$\square$

\remarkheading{Three uses of a reciprocal energy}{rem:gen-appended-energy}
Let $U$ be a measurable energy on $\mathsf Y$, bounded below, that couples coordinates reciprocally. If $U$
enters a potential \eqref{eq:gen-gibbs-potential} with $Z_X=+\infty$, no law is defined and there is no
evidence to bound. If it enters with $0<Z_X<\infty$, the undirected model is a legitimate joint to which
\Cref{ch:elbo} applies. If $U$ is appended to the free energy of a joint $\mathbb P_{\theta,V}$ that does not
contain it, the evidence of that joint is unaffected, and for every recognition law $Q$ in the classical
domain of \Cref{hyp:elbo-evidence-domain} with $\E_Q\lvert U\rvert<\infty$,
\begin{equation}
\Fenergy_V[Q;X,o]+\E_QU=\widetilde\Fenergy_V[Q;X,o]-\log Z_U,
\qquad
Z_U=\int e^{-U(y)}\mathbb P_{\theta,V}(dy,do\given X)\in(0,\infty),
\label{eq:gen-appended-energy}
\end{equation}
where $\widetilde\Fenergy_V$ is the free energy of the tilted joint
$Z_U^{-1}e^{-U(y)}\mathbb P_{\theta,V}(dy,do\given X)$. The appended objective is therefore an evidence bound of
that tilted joint only after $\log Z_U$, which depends on $(\theta,X)$, is restored. \status{ESTABLISHED}

\noindent\emph{Proof.}
Since $U$ is bounded below, $0<Z_U\leq\sup e^{-U}<\infty$, and the tilted joint has density
$p_{\theta,V}(o,y\given X)e^{-U(y)}/Z_U$. With $\Fenergy_V[Q;X,o]=\E_Q[\log q_{V,o,X}(Y)-\log p_{\theta,V}(o,Y\given X)]$,
adding $\E_QU$ and splitting $\log Z_U$ off the tilted log density gives \eqref{eq:gen-appended-energy}; all
expectations are finite under the stated integrability. \hfill$\square$

The directed construction is used because its normalization does not depend on the parameters.

\section{Covariance under passive frame re-expression}
\label{sec:gen-gauge}

A passive re-expression $(g_i)_{i\in V}$ acts on samples through $R_i^x=\rho_x(g_i)$ of
\eqref{eq:geo-represented-frame-change}, and on the latent and joint spaces by the first two lines
below, so records are gauge inert. The structural data transform to $X'$ with the frames
$u_i\cdot g_i^{-1}$ and, by the vertex law \eqref{eq:geo-regime-two-gauge-law}, the links of the third
line; a structural constraint tying parallel copies $e,e':j\to i$ to one link is preserved, because both
copies transform by the same pair (fourth line). Active invariance is invariance under the one copy
$\mathcal A_{c_*}\cong G$ of \Cref{prop:geo-active-image-point}.
\begin{gather}
R=\bigoplus_{i\in V}\bigl(R_i^b\oplus R_i^m\bigr):\mathsf Y\longrightarrow\mathsf Y,
\label{eq:gen-rep-actions}\\
T(y,o)=(Ry,o),
\label{eq:gen-contextual-frame-change}\\
\Theta_e^{x\prime}=g_i\Theta_e^xg_j^{-1},
\qquad e:j\to i,\ x\in\{b,m\},
\label{eq:gen-gauge-links}\\
\Theta_e^x=\Theta_{e'}^x
\quad\Longrightarrow\quad
g_i\Theta_e^xg_j^{-1}=g_i\Theta_{e'}^xg_j^{-1}.
\label{eq:gen-shared-link-admissibility}
\end{gather}

\hypothesisheading{Kernel covariance}{hyp:gen-kernel-covariance}
The transformed parameters $(\theta',X')$ are such that every factor obeys the first two lines below,
with the empty parent tuple at roots, and every record kernel obeys the third, where $R_{\partial a}$ is
the restriction of $R$ to the incident blocks; for a record with scope $\{i\}$ the third line reads as
the fourth:
\begin{gather}
P_{\theta',i}^{m}\bigl(\cdot\given R_{\operatorname{pa}(i)}^mm_{\operatorname{pa}(i)},X'\bigr)
=(R_i^m)_\#P_{\theta,i}^{m}\bigl(\cdot\given m_{\operatorname{pa}(i)},X\bigr),
\label{eq:gen-gauge-pushforward-model}\\
P_{\theta',i}^{k}\bigl(\cdot\given R_{\operatorname{pa}(i)}^bk_{\operatorname{pa}(i)},R_i^mm_i,X'\bigr)
=(R_i^b)_\#P_{\theta,i}^{k}\bigl(\cdot\given k_{\operatorname{pa}(i)},m_i,X\bigr),
\label{eq:gen-gauge-pushforward-state}\\
K_a\bigl(X',R_{\partial a}y_{\partial a},\cdot\bigr)=K_a\bigl(X,y_{\partial a},\cdot\bigr),
\qquad a\in\mathfrak A,
\label{eq:gen-gauge-record-kernel}\\
K_a\bigl(X',(R_i^bk_i,R_i^mm_i),\cdot\bigr)=K_a\bigl(X,(k_i,m_i),\cdot\bigr).
\label{eq:gen-gauge-pushforward-obs}
\end{gather}
The hypothesis excludes mechanisms whose coordinate form is compatible with no re-expression of the
parameters.

\propositionheading{Joint measure covariance and evidence invariance}{prop:gen-product-evidence-invariance}
Under \Cref{hyp:gen-kernel-covariance}, with $R$ and $T$ as in
\eqref{eq:gen-rep-actions}--\eqref{eq:gen-contextual-frame-change}, the first line below holds as
probability measures, without a density or reference hypothesis, and the observation marginals agree, so
the second holds:
\begin{gather}
\mathbb P_{\theta',V}(\cdot\given X')=T_\#\mathbb P_{\theta,V}(\cdot\given X),
\label{eq:gen-gauge-joint-pushforward}\\
p_{\theta',V}(o\given X')=p_{\theta,V}(o\given X)
\quad\text{for }\nu^{\mathsf O}\text{-almost every }o.
\label{eq:gen-gauge-evidence-invariance}
\end{gather}
If latent densities exist and $R_\#\nu^{\mathsf Y}\ll\nu^{\mathsf Y}$, then with $j_R$ of the first line
below the second holds $(\nu^{\mathsf Y}\otimes\nu^{\mathsf O})$-almost everywhere, so composing a density
with $R^{-1}$ represents the pushforward for every density exactly when $j_R=1$ almost everywhere. If every
represented block is an invertible linear map between Euclidean spaces of equal dimension, with Lebesgue
reference (equal within each channel, with $\dim\mathsf K\neq\dim\mathsf M$ allowed), then
$j_R=\lvert\det R\rvert^{-1}$ and the third line holds:
\begin{gather}
j_R(y'):=\bigl(d(R_\#\nu^{\mathsf Y})/d\nu^{\mathsf Y}\bigr)(y'),
\label{eq:gen-gauge-reference-rn}\\
p_{\theta',V}(o,y'\given X')=p_{\theta,V}(o,R^{-1}y'\given X)j_R(y'),
\label{eq:gen-gauge-rn-density}\\
p_{\theta',V}(o,y'\given X')=\lvert\det R\rvert^{-1}p_{\theta,V}(o,R^{-1}y'\given X).
\label{eq:gen-gauge-lebesgue-density}
\end{gather}
Finally, the root law bridges \eqref{eq:gen-root-law-bridge} of $\theta$ and $\theta'$ satisfy
$\widetilde\Phi^{\mathrm{gen}\prime}_i=\widehat\rho_b(g_i)\circ\widetilde\Phi^{\mathrm{gen}}_i\circ\widehat\rho_m(g_i)^{-1}$,
the law \eqref{eq:geo-tildephi-gauge-law} of a cross-morphism representative. \status{ESTABLISHED}

\noindent\emph{Proof.}
Apply \eqref{eq:gen-gauge-pushforward-model} and \eqref{eq:gen-gauge-pushforward-state} along the
topological ordering: finite kernel induction gives
$\mathbb P_{\theta',V}^{\mathsf Y}(\cdot\given X')=R_\#\mathbb P_{\theta,V}^{\mathsf Y}(\cdot\given X)$. For bounded
measurable $F$ on $\mathsf Y\times\mathsf O$, write $\mathbb P_{\theta',V}(\cdot\given X')$ by
\eqref{eq:gen-record-extension}, substitute $y'=Ry$ in the latent integral, and replace the record kernels
at $(X',R_{\partial a}y_{\partial a})$ by those at $(X,y_{\partial a})$ by \eqref{eq:gen-gauge-record-kernel};
this gives $\int F(y,o)\mathbb P_{\theta',V}(dy,do\given X')=\int F(T(y,o))\mathbb P_{\theta,V}(dy,do\given X)$,
which is \eqref{eq:gen-gauge-joint-pushforward}, and projecting to $o$ gives
\eqref{eq:gen-gauge-evidence-invariance}.
For $Q=q\nu^{\mathsf Y}$ and measurable $B$,
$R_\#Q(B)=\int_{R^{-1}B}q(y)\nu^{\mathsf Y}(dy)=\int_Bq(R^{-1}y')(R_\#\nu^{\mathsf Y})(dy')$, so
$R_\#Q=(q\circ R^{-1})j_R\nu^{\mathsf Y}$; applied to the joint density this is
\eqref{eq:gen-gauge-rn-density}, and linear change of variables gives $j_R=\lvert\det R\rvert^{-1}$. For the
bridge, \eqref{eq:gen-gauge-pushforward-state} at a root gives
$\widetilde\Phi^{\mathrm{gen}\prime}_i(\widehat\rho_m(g_i)s)=\int P_{\theta',i}^{k}(\cdot\given R_i^mm,X')s(dm)
=\widehat\rho_b(g_i)\widetilde\Phi^{\mathrm{gen}}_i(s)$. \hfill$\square$

\corollaryheading{Independent channel-frame invariance}{cor:elbo-independent-frame-invariance}
Under \Cref{hyp:gen-kernel-covariance}, let $\mathbb Q_{V,o,X}$ be selected before the re-expression and
transported as a whole, $\mathbb Q'_{V,o,X'}=R_\#\mathbb Q_{V,o,X}$. Then for $\mathbb P_{\theta,V,X}^{\mathsf O}$-almost
every $o$ the posterior transforms as $\boldsymbol\Pi_{\theta',V,o,X'}=R_\#\boldsymbol\Pi_{\theta,V,o,X}$, and
\begin{equation}
\KL(\mathbb Q'_{V,o,X'}\Vert\boldsymbol\Pi_{\theta',V,o,X'})
=\KL(\mathbb Q_{V,o,X}\Vert\boldsymbol\Pi_{\theta,V,o,X}),
\quad
\Lelbo'^{\mathrm{ext}}_V(\mathbb Q'_{V,o,X'};X',o)=\Lelbo^{\mathrm{ext}}_V(\mathbb Q_{V,o,X};X,o)
\label{eq:elbo-gauge-invariance}
\end{equation}
in the extended reals, the singular branch included, with $\Lelbo^{\mathrm{ext}}_V$ as in
\eqref{eq:elbo-extended-definition}; on the classical domain of both sides the same holds for $\Lelbo_V$.
\status{ESTABLISHED}

\noindent\emph{Proof.}
By \eqref{eq:gen-gauge-joint-pushforward}, for measurable $B\subseteq\mathsf Y$ and $A\subseteq\mathsf O$,
$\mathbb P_{\theta',V}(B\times A\given X')=\mathbb P_{\theta,V}(R^{-1}B\times A\given X)
=\int_A(R_\#\boldsymbol\Pi_{\theta,V,o,X})(B)\mathbb P^{\mathsf O}_{\theta,V,X}(do)$, and the observation marginals
agree, so $R_\#\boldsymbol\Pi_{\theta,V,o,X}$ is a version of the new posterior. \Cref{lem:prob-kl-chain-rule}(e)
gives equality of the divergences in $[0,+\infty]$. The evidence agrees by
\eqref{eq:gen-gauge-evidence-invariance}, and substitution in \eqref{eq:elbo-extended-definition} subtracts
equal quantities from equal finite numbers; the classical functional equals the extended one on its domain.
\hfill$\square$

The running example satisfies \Cref{hyp:gen-kernel-covariance} with $\theta'=\theta$, since $\rho_m$ is
trivial and its isotropic state covariances, link-transported means, and record lengths are compatible
with \eqref{eq:gen-gauge-links}.

\section*{Physicist's summary}

The generative model is one normalized joint of models, beliefs given parents, and records, which may
close arbitrary cycles; reciprocal latent conditionals define no joint. No factor may read the
recognition law it is scored against, and passive frame changes preserve evidence and bound.
